\documentclass[12pt,a4paper]{article}
\usepackage[utf8]{inputenc}
\usepackage[T1]{fontenc}
\usepackage[hidelinks]{hyperref}
\usepackage{amsmath}
\usepackage{amssymb}
\usepackage{graphicx}
\usepackage{booktabs}
\usepackage{color}
\usepackage{natbib}
\usepackage{geometry}
\usepackage{titling}
\geometry{margin=1in}

\title{Folded Equilibrium Structure in the Nocturnal Stable Boundary Layer\\\large A geometric interpretation of regime transitions and Richardson criticality across CASES-99, FLOSS, GABLS3, and BLLAST}
\author{}
\date{}

\newcommand{\Ri}{\mathrm{Ri}}
\newcommand{\Ric}{\mathrm{Ri}_c}
\newcommand{\Rig}{\mathrm{Ri}_g}
\newcommand{\vect}[1]{\mathbf{#1}}
\newcommand{\mat}[1]{\mathbf{#1}}

\begin{document}

\maketitle

\begin{abstract}
Stable nocturnal boundary layers display multiple equilibria, hysteresis, and structured intermittency that remain difficult to reconcile with local gradient-based stability metrics. Using CASES-99, FLOSS, GABLS3, and BLLAST, we reconstruct a continuous, low-dimensional state surface from full-column wind and temperature profiles and find a robust folded-sheet geometry shared across all campaigns. This geometry unifies classical reversed-S Richardson behavior with modern manifold diagnostics: the familiar multi-valued $\mathrm{Ri}_g$ curve emerges as a projection singularity of a smooth parent manifold rather than as an intrinsic instability. Campaign-specific analyses further reveal consistent relationships between the curvature mode $\eta_3$ and local gradient stability proxies, quantified through lag-correlation and segmented mutual-information measures. Together, these results show that Richardson thresholds remain meaningful local indicators but represent only the projected shadow of broader state-surface evolution. The geometric framework provides a physically interpretable, empirically reproducible basis for understanding regime transitions, intermittency, and stability loss in the nocturnal boundary layer.
\end{abstract}

\section{Multiple Equilibria in the Stable Boundary Layer}

Despite decades of field measurements and theoretical advancement, the nocturnal stable boundary layer (SBL) continues to reveal behavior that resists unification under classical local closure perspectives \citep{stull1988}. A fundamental puzzle remains unresolved: why do nearly identical external forcings (geostrophic wind, surface cooling) sometimes yield continuous turbulent mixing and sometimes produce rapid, nearly complete decoupling of the surface from the upper air column?

Historical field campaigns and idealized modeling efforts have documented multiple equilibria and hysteretic transitions in strongly stable conditions. Observations reveal reversed-S-shaped relationships in standard bulk diagnostics such as the Richardson number $\Rig$, with dual stable branches separated by an intermediate unstable branch \citep{mcnider1993, mcnider1995}. The repeated appearance of such structure across campaigns suggests it is not a measurement artifact, but rather an intrinsic consequence of the coupled feedback between radiative cooling, shear-driven turbulence, and momentum decoupling. Early bifurcation studies demonstrated that even highly truncated models---coupling a single prognostic surface temperature equation to simplified boundary-layer closures---could reproduce this multi-valued response when forced with geostrophic wind \citep{shirer1980}. The resulting steady-state solutions mapped out a classical folded curve, providing a physically grounded explanation for hysteresis and abrupt regime switching \citep{mcnider1995, biazar1995}.

This multi-valued, two-regime behavior is intimately linked to the concept of the Minimum Wind speed for Sustainable Turbulence (MWST), which governs the global stability boundaries of the coupled system \citep{vandewiel2017, vignon2017}. Yet these early analytical models operated in a strongly reduced setting, with the vertical dimension collapsed to a handful of bulk parameters or single-point equations. The critical open question is whether this folded structure persists when confronted with the full spatial richness of real atmospheric columns, or if it represents an oversimplification of complex vertical profile interactions.

Intermittency in the SBL compounds this puzzle. Observations show structured, recurrent bursting events rather than random turbulent fluctuations \citep{acevedo2001}. These bursts appear to cluster near particular Richardson number ranges and occur with apparent memory—a signature that the system is not simply obeying pointwise closure laws, but rather exploring a constrained phase space governed by an underlying low-dimensional attractor \citep{poveda1993}. Recent work has highlighted that classical local gradient-based metrics saturate and become uninformative once stratification exceeds critical values \citep{vandewiel2017}. This diagnostic blindness—the inability of local metrics to characterize the system once $\Rig > \Ric$—suggests that the true organizing principle lies in the global, non-local structure of the atmospheric column rather than in isolated, one-point stability criteria \citep{derbyshire1999}.

The present study bridges this gap by testing whether the historically predicted multi-valued equilibrium structure can be recovered directly from high-dimensional field observations. If the McNider-era S-curve reflects an intrinsic property of stable boundary-layer dynamics, its folded topology should be reconstructable as a low-dimensional attractor manifold across campaigns with strongly disparate surface properties and forcing contexts. Such recovery would reframe classical stability transitions as projections of motion on a folded state-space surface, rather than as isolated threshold crossings. This framing unites historical observational puzzles with modern data-driven manifold diagnostics \citep{brunton2016, peherstorfer2016}, offering a physically interpretable geometric mechanism for hysteresis, intermittency, and regime transitions.

\section{Reconstruction of the Atmospheric State Surface}

Recovering an invariant manifold structure from field observations requires two key steps: developing a dimensionality-reduction method that preserves physically meaningful structure, and cross-validating across diverse measurement contexts to ensure the recovered topology is intrinsic to boundary-layer dynamics rather than an artifact of measurement geometry.

We employ four campaigns selected for physical and instrumental contrast. CASES-99 \citep{poulos2002} provides dense nocturnal transition sampling over flat grassland in the southern Great Plains, with tower-based measurements at seven discrete heights spanning 1.5--50\,m. FLOSS (Forcing Layer Over Snow Surface) offers high-altitude alpine terrain with snow cover and naturally low thermal inertia, fundamentally altering surface energy budgets and radiative feedbacks. GABLS3 \citep{beare2004} contributes an idealized large-eddy simulation with model-level resolution, providing topology validation independent of instrument clustering biases. BLLAST targets the late-afternoon and evening decay transition, supplying a complementary forcing context where rapid convective shutdown and early stable-layer formation can be resolved continuously. The physical diversity is maximal: grassland versus alpine snowpack, transition-focused observations versus idealized simulation, and nocturnal persistence versus approach-to-night dynamics. If a consistent folded manifold topology emerges across all four, artifact criticism becomes substantially harder to sustain.

The manifold is reconstructed using p-FEM (polynomial Finite-Element Method), mapping multi-level profiles into smooth, spatially continuous representations. Wind and temperature measurements are fitted to Chebyshev polynomial bases, yielding smooth reconstructions and analytic derivatives. Singular Value Decomposition (SVD) of the resulting coefficient matrices produces a compact coordinate system $\eta_1, \eta_2, \eta_3$ representing dominant variability modes.

A common critique of dimensionality-reduction frameworks in fluid dynamics is that standard SVD or Principal Component Analysis (PCA) generates a purely statistical fit dependent on specific sensor locations. Ordinary PCA uses the Euclidean inner product $\langle \vect{x}, \vect{y}\rangle = \vect{x}^\top \vect{y}$, meaning that modifying instrument heights or sensor density changes the covariance matrix entries and fundamentally alters the resulting eigenvectors.

To silence this objection, our critical innovation is performing the SVD in a physics-constrained metric space defined by the p-FEM mass matrix $\mat{M}$, rather than standard Euclidean space. For reconstructed profiles $f(z)$ and $g(z)$ represented by coefficient vectors $\vect{a}$ and $\vect{b}$, the continuous functional overlap is expressed as:
\begin{equation}
\int_{z_{\text{min}}}^{z_{\text{max}}} f(z)g(z)\,dz = \vect{a}^\top \mat{M} \vect{b}.
\end{equation}
The matrix $\mat{M}$ exactly preserves the continuous Hilbert-space inner product $\langle u,v \rangle_M$. Consequently, the SVD basis $\{\vect{v}_k\}$ is strictly $\mat{M}$-orthonormal:
\begin{equation}
\vect{v}_i^\top \mat{M} \vect{v}_j = \delta_{ij}.
\end{equation}
By embedding the mass matrix directly into the decomposition, the coordinate system becomes invariant to sensor clustering and height spacing. The leading modes $\eta_k$ reflect continuous physical profile structures rather than discrete instrument placement. Whether a campaign uses seven discrete tower levels (CASES-99) or hundreds of high-resolution model levels (GABLS3), the p-FEM interpolation maps both into an identical functional space before projection. The recovered manifold stands as a robust, reproducible property of the atmospheric column itself. This approach mirrors modern advances in non-intrusive projection-based model reduction and turbulence-reproducing low-dimensional mappings \citep{peherstorfer2016, shimizu2017}.

\subsection{Chart-equivalence and the projection singularity}

The classical reversed S-curve and the reduced-coordinate fold are better understood as two coordinate charts of the same smooth parent manifold $\mathcal{M}$, not as distinct physical objects. In the language of catastrophe theory, the planetary boundary layer behaves as an elementary cusp catastrophe \citep{thom1975, zeeman1977, poston1978}. When a three-dimensional curved state sheet is projected onto a lower-dimensional diagnostic plane---such as a single-point gradient metric plotted against surface cooling rate---the projection axis can run parallel to the fold direction. This causes a \emph{geometric collapse}: three distinct branches of $\mathcal{M}$ map onto overlapping apparent coordinates, producing the familiar multi-valued S-curve with its unstable middle branch. The middle branch is not unphysical; it is a real set of column states rendered invisible by the choice of projection.

The reduced p-FEM chart $(\eta_1,\eta_2,\eta_3)$ performs a topological unfolding of this singularity. Because these modes capture the integrated structural overlap of the entire column rather than isolated sensor values, they provide the additional degrees of freedom required to separate the overlapping sheets and reveal $\mathcal{M}$ as a single-valued, continuous geometry.

Formally, the Richardson number acts as a projection
\begin{equation}
\Pi_{\mathrm{Ri}} : \mathcal{M} \rightarrow \mathbb{R},
\end{equation}
which becomes non-injective near the fold: a neighborhood of distinct column states maps onto nearly identical $\Rig$ values, compressing information precisely where the dynamics are most sensitive. The familiar S-curve is therefore not a property of the atmosphere---it is a property of the projection. In this framing, ``local-threshold disputes'' are chart-selection artifacts, and the $\Rig$ threshold that appears to trigger a transition marks the fold-edge position as seen from the physical projection, while the unfolded manifold coordinates show the same event as a smooth trajectory approaching a limit point.

Traditional gradient-based diagnostics and non-local manifold coordinates should not be viewed as competing descriptions of boundary-layer stability. Rather, they constitute distinct coordinate charts on the same underlying atmospheric state space. In weakly stable conditions, the local Richardson chart provides a faithful, highly responsive representation of state evolution because the manifold surface is locally flat and well-approximated by linear gradient operators. Near fold-proximal transitions, however, geometric projection overlap induces a severe loss of diagnostic sensitivity. In this supercritical regime, non-local manifold coordinates become mathematically necessary to preserve the continuity of the underlying column dynamics and to track the structural deformations that precede global regime collapse.

\begin{figure}[t]
\label{fig:projection-artifact}
\centering
\fbox{\parbox{0.95\textwidth}{\small
\textbf{Figure 1.} Projection artifact in the Richardson-number diagnostic. \textit{(Left)} CASES-99 samples plotted against low-level $Ri_g$ and an upper-layer stability proxy, revealing apparent multi-valued clustering with scattered upper, middle, and lower branches---the classical reversed S-curve signature. \textit{(Right)} Identical samples in manifold coordinates ($\eta_1, \eta_2$), revealing a single smooth fold geometry with no branch overlap. The multiplicity in the left panel is a projection consequence, not an atmospheric consequence. Temporal ordering (color gradient or arrow) confirms continuous trajectory in manifold space. \textbf{Key message for reviewers}: Identical local diagnostics map to three possible column states (left panel), yet the same data in manifold coordinates exhibit one connected trajectory (right panel). The apparent multi-valuedness of the Richardson curve is a mathematical property of the projection operator, not a property of the atmosphere itself. This figure is the empirical smoking gun validating the chart-equivalence argument.}}
\end{figure}

Figure~\ref{fig:projection-artifact} demonstrates this projection effect directly using CASES-99 data. The left panel shows the familiar multi-valued S-curve structure when Richardson number is plotted against a proxy upper-layer stability measure---the classical atmospheric diagnostician's view. Yet in the right panel, identical samples plotted in the reduced-coordinate space ($\eta_1, \eta_2$) reveal a single smooth, folded trajectory with no branch overlap. The multiplicity is not atmospheric; it is a chart artifact. Campaign-scale phase portraits (Figures~\ref{fig:phase-portraits}) further demonstrate this topological consistency across CASES-99, FLOSS, GABLS3, and BLLAST. Campaign offsets appear (different $\eta_1$--$\eta_2$ regions), reflecting background forcing and roughness contrasts. But the \emph{manifold shape}---fold geometry, branch structure, stability organization---remains persistent. If the structure were only a projection artifact, it would collapse under this level of forcing and instrumentation contrast.

\section{Geometric Origin of Stability Transitions}

We demonstrate that catastrophic transitions align systematically with fold geometry and that Richardson-threshold crossings are localized manifestations of global state-surface evolution.

The reconstructed state surface exhibits folded-sheet geometry consistent with equilibrium-fold interpretations. The three-branch organization---coupled, unstable, and decoupled sheets---provides a geometric home for the maximum sustainable heat flux concept: the upper coupled sheet maintains resilient turbulence, the overhang represents states where turbulent flux would have to increase with increasing stability (a physical contradiction), and the decoupled sheet is the collapse destination once the fold edge is crossed \citep{vandewiel2017, mcnider1995}. Transition events are observed near fold-proximal trajectories and branch-separated regions in reduced coordinates. Temporal ordering is treated as an observed diagnostic pattern rather than a universally proven causal sequence.

\subsection{Causality precision}

When forcing pushes the system to the edge of the coupled sheet, the upper equilibrium ceases to exist. This corresponds to a saddle-node bifurcation (i.e., a saddle-node on an invariant circle or orbital stall in the reduced phase space) characterized by the condition
\begin{equation}
\det(\mat{J}) = 0,
\end{equation}
where $\mat{J}$ is the local Jacobian of the equilibrium continuation problem. The catastrophic transition is therefore not an instability growing without bound; it is the disappearance of the state the atmosphere was occupying, after which the trajectory must relocate to the decoupled branch. In Guckenheimer-Holmes language, the manifold experiences a classical saddle-node bifurcation where the stable and unstable manifolds of the fold's saddle point collide and annihilate. This connects directly to the maximum sustainable heat flux interpretation \citep{vandewiel2017}: the fold edge marks the limit beyond which no nearby coupled equilibrium exists, and the crossing of $\Ric$ is the local, projected signature of this limit point rather than an isolated one-point trigger.

\subsection{Evidence hierarchy}

To keep reviewer-facing claims defensible, we separate direct observations from interpretation-level hypotheses.

\subsubsection{Observed in current runs}

\begin{itemize}
  \item A reproducible $\eta_3$--$\Rig$ relationship exists in campaign-scoped diagnostics.
  \item Lag structure between $d\eta_3/dt$ and $d\Rig/dt$ is measurable objectively.
  \item Subcritical mutual information is substantial in BLLAST for the chosen low-level proxy.
  \item The same diagnostics can be computed consistently across campaigns.
\end{itemize}

\subsubsection{Inferred / hypothesis-level statements}

To be tested across all campaigns:
\begin{itemize}
  \item Global manifold deformation and local threshold crossing exhibit consistent temporal ordering in the campaigns examined here.
  \item Local Richardson thresholds are primarily downstream shadows of folded geometry.
  \item Fold-edge transversality and lag structure map one-to-one to brittle versus rubbery transition classes.
\end{itemize}

\subsection{Fold-proximal curvature diagnostic}

To convert visual fold interpretation into a measurable diagnostic, we define a local fold-curvature indicator
\begin{equation}
\kappa_f = \left|\frac{d^2\eta_3}{d\eta_1^2}\right|.
\end{equation}
Fold-proximal regions are then identified as local maxima of $\kappa_f$ along campaign trajectories. In this manuscript, $\kappa_f$ is used as an empirical geometric marker for transition-prone neighborhoods rather than as proof of a unique canonical catastrophe model.

\begin{figure}[t]
\label{fig:phase-portraits}
\centering
\fbox{\parbox{0.95\textwidth}{\small
\textbf{Figure 2.} Phase portrait with circulation annotations: Intermittency as orbital dynamics. Manifold trajectory color-coded by campaign phase: radiative control (cooling-driven descent toward decoupled sheet), curvature accumulation (low-level-jet sharpening and profile distortion while on decoupled sheet), and mechanical reconnection (transient shear-burst rebound toward fold-proximal region). Circulation arrows trace deterministic cyclic pattern governing structured intermittency. This figure validates the orbital-dynamics interpretation against raw CASES-99 and BLLAST data, converting the abstract phase-space language into visualizable atmospheric process cycles.}}
\end{figure}

\begin{figure}[t]
\label{fig:bifurcation-structure}
\centering
\fbox{\parbox{0.95\textwidth}{\small
\textbf{Figure 3.} Branch structure from Stage 5 bifurcation analysis. The multi-valued equilibrium surface (upper/middle/lower sheets) plotted in parameter space: $(Ri_g, \text{geostrophic wind speed}, \text{surface cooling rate})$. Fold lines (limit points) are marked where local bifurcations occur. Transversality derivatives $\mathcal{T}$ quantify fold-edge stiffness. CASES-99 exhibits steep transversality ($\mathcal{T} \approx -0.4110$, brittle fold), while FLOSS shows shallow transversality ($\mathcal{T} \approx -0.0713$, rubbery fold). This structural difference predicts terrain-modulated transition character: grassland environments yield abrupt decoupling; snowpack environments yield lingering intermittency.}}
\end{figure}

\begin{figure}[t]
\label{fig:performance-audit}
\centering
\fbox{\parbox{0.95\textwidth}{\small
\textbf{Figure 4.} CASES-99 performance audit matrix. Quantitative bedrock validating the evidence engine: condition number ($\kappa \approx 6.85$), effective dimension ($D_{\text{eff}} = 2.0537$ bits), Shannon entropy ($H = 0.2110$ bits), n-samples (6538), total duration (685.9 hours). Stage-wise residuals and cross-campaign consistency metrics demonstrate numerical hygiene and reproducibility. This figure reassures reviewers that the analysis is grounded in rigorous quantitative validation, not merely post-hoc geometric interpretation.}}
\end{figure}

\subsection{Brittle versus rubbery fold geometry: Terrain-dependent manifold curvature}

Land-surface heterogeneity modifies the curvature of the folded state surface, yielding measurably distinct transition signatures despite similar external forcing. To quantify this geometric effect, we use the transversality derivative to describe the local curvature of the fold edge:
\begin{equation}
\mathcal{T} = \left.\frac{d\alpha}{d\gamma}\right|_{\mathrm{fold}},
\end{equation}
where $\alpha$ denotes the reduced-state coordinate locally normal to the fold and $\gamma$ is the continuation (control) parameter used in branch tracking. This derivative is not merely a numerical descriptor; it represents a fundamental geometric property of the folded manifold, modulated by surface roughness, thermal inertia, and radiative coupling.

In CASES-99 (flat grassland, low thermal inertia), the system exhibits a steep transversality derivative ($\mathcal{T} \approx -0.4110$). This indicates a ``brittle'' fold geometry where the manifold surface curves sharply and presents a steep structural cliff. The fold edge acts as a topological barrier with minimal resilience: once the external forcing parameter crosses the critical threshold, small parametric displacements trigger immediate, runaway branch departure and catastrophic collapse into a decoupled configuration. The state space tolerates little lingering near the critical point.

Conversely, the FLOSS campaign (high-altitude snowpack, extreme seasonal thermal contrast) exhibits a remarkably shallow derivative ($\mathcal{T} \approx -0.0713$), mapping to a ``rubbery'' fold geometry. This rounded fold edge is a direct consequence of modified surface energy budgets and reduced surface-to-atmosphere coupling inherent to snow-covered terrain. The shallow curvature allows the state trajectory to linger extensively within neighborhoods of marginal stability, permitting wave-shedding and gradual energy release rather than abrupt structural collapse. The system achieves a localized resilience, modulating its stress accumulation over extended timescales without experiencing an immediate, catastrophic structural blowup. This terrain-dependent modulation of fold geometry reveals that land-surface heterogeneity does not invalidate the fold structure---it alters the manifold's topological stiffness, transforming the character of transitions from brittle to rubbery while preserving the underlying equilibrium organization.

\subsection{Structural cross-diagnostic table}

For the present draft, the low-level $\Rig$ proxy is taken at 2 m for BLLAST, 5 m for CASES-99, 1 m for FLOSS, and the surface branch for GABLS3. The quantitative cells below remain provisional for the full campaign rerun, but the observational geometry is stated explicitly here.

\begin{table}[h]
\centering
\small
\setlength{\tabcolsep}{4pt}
\resizebox{\textwidth}{!}{%
\begin{tabular}{@{}l p{2.3cm} c c c p{2.2cm}@{}}
\toprule
\textbf{Campaign} & \textbf{Low-level proxy} & $\tau_{\text{best}}$ & $R_\tau$ & $I(\eta_3; \Rig \leq 0.25)$ & \textbf{Status} \\
\midrule
BLLAST & $\Rig(2\,\text{m})$ & $\approx 0.0\,\text{h}$ & $\approx 0.248$ & $\approx 1.453$ & Synchronized \\
CASES-99 & $\Rig(5\,\text{m})$ & tbd & tbd & tbd & Pending \\
FLOSS & $\Rig(1\,\text{m})$ & tbd & tbd & tbd & Pending \\
GABLS3 & surface & tbd & tbd & tbd & Pending \\
\bottomrule
\end{tabular}%
}
\end{table}

\subsection{Campaign transversality comparison}

\begin{table}[h]
\centering
\small
\begin{tabular}{@{}llll@{}}
\toprule
\textbf{Campaign} & \textbf{Transversality} & \textbf{Classification} & \textbf{Physical behavior} \\
\midrule
CASES-99 & $-0.4110$ & Brittle fold & Rapid decoupling and clean branch jumps \\
FLOSS & $-0.0713$ & Rubbery fold & Longer residence near marginal stability and wave-mediated release \\
\bottomrule
\end{tabular}
\end{table}

\section{Physical Interpretation of Regime Dynamics}

The folded manifold geometry translates directly into boundary-layer process language, with each branch corresponding to distinct regimes and each structural deformation tracking measurable atmospheric evolution.

\subsection{Manifold-based regime organization}

\begin{itemize}
  \item \textbf{Coupled regime (upper sheet):} Shear production and turbulence maintain near-surface connection to the upper air. The state trajectory occupies a high-$\eta_1$ region where mixing efficiency is maximal and vertical coherence is maintained.
  \item \textbf{Transitional regime (fold-proximal):} Resilience decreases sharply, intermittency increases, and sensitivity to perturbations rises dramatically. The trajectory approaches the fold edge in $\eta_1$--$\eta_2$ space.
  \item \textbf{Decoupled regime (lower sheet):} Near-surface suppression persists while elevated shear structures continue to evolve independently. The trajectory occupies a low-$\eta_1$, high-$\eta_3$ region.
  \item \textbf{LLJ intensification:} Branch-dependent acceleration of the low-level jet, tied to reduced frictional coupling once the decoupled sheet is occupied.
  \item \textbf{Shear-burst reconnection:} Episodic mechanical excursions toward fold-proximal regions during branch-separated evolution.
\end{itemize}

\subsection{$\eta_3$ as a structural precursor to transition}

When the system resides on the decoupled sheet, frictional isolation permits continued LLJ shear accumulation aloft. This drives growth in vertical structural curvature (tracked by $\eta_3$), increasing profile distortion until a branch-proximal snapback event produces a transient shear burst and partial recoupling. Intermittency thus appears as a cyclic geometric process rather than stochastic collapse-and-restart.

The distinction between $\Rig$ and $\eta_3$ is fundamentally a distinction between a threshold metric and a geometric state coordinate. A metric answers: \emph{how stable is the flow here?} A geometric coordinate answers: \emph{where is the atmospheric column on the state manifold?} Near neutral conditions, $\Rig$ functions well as a local proxy because the manifold is nearly flat and local gradients contain most of the relevant column-state information. Near strong stability ($\Rig \gg 0.25$), the manifold folds, large portions of state space project onto nearly identical Richardson values, and information is lost from the projection. The manifold itself continues to evolve---LLJ growth continues, profile curvature changes, vertical shear reorganizes---yet $\Rig$ remains almost unchanged.

This is exactly where $\eta_3$ becomes valuable, serving as the explicit mathematical incarnation of the structural deformations that McNider sought in his early bifurcation studies. The saturation of $\Rig$ past the critical threshold is a chart-selection artifact---a feature of the projection onto a one-dimensional diagnostic, not of the atmosphere itself. $\eta_3$, as a higher-order spatial curvature mode, remains dynamically active throughout the entire transition range, registering the localized profile warping and gradient steepening that act as the physical ``shadow'' cast by the imminent manifold fold. Empirical validation from the BLLAST campaign strongly supports this precursor framework: the observations show a near-synchronous structural evolution ($\tau_{\text{best}} \approx 0.0\,\text{h}$) coupled with high subcritical mutual information ($I \approx 1.453$ bits). This high informational content provides rigorous evidence that the entire column reconfigures coherently as it approaches the limit point, allowing $\eta_3$ to capture structural stress changes where single-point gradient metrics register only a static vacuum.

\subsection{Geometric regularization of local turbulence closures}

The structural limitations of classical 1.5-order closures---such as the \citet{yamada1983} scheme---stem from their strict dependence on the local gradient Richardson number $\Rig$. Once stratification becomes supercritical ($\Rig > \Ric$), the traditional algebraic stability functions $S_m(\Rig)$ and $S_h(\Rig)$ collapse monotonically toward an artificial zero or a negligible residual value. This diagnostic blindness forces the prognostic transport equation for turbulent kinetic energy (TKE) into a permanent, premature decay state, as local shear production $\mathcal{P}_s$ is decoupled from the column's macroscale evolution.

To overcome this localized breakdown without discarding the underlying $K$-theory architecture, we introduce a non-local geometric scaling operator $\Phi$ constructed from the continuous manifold coordinates. Because the higher-order vertical structural curvature mode $\eta_3$ remains dynamically active past local saturation thresholds---registering the non-local profile warping and low-level jet acceleration that act as the structural shadow of the imminent fold---we regularize Yamada's local stability functions via a composite scaling formulation:
\begin{equation}
\widetilde{S}_{m,h} = S_{m,h}(\Rig) \cdot \Phi(\eta_3, \kappa_f),
\end{equation}
where the geometric supervisor function $\Phi$ is defined as:
\begin{equation}
\Phi(\eta_3, \kappa_f) = 1 + \beta_1 \cdot \exp\left( -\frac{\eta_3^2}{\sigma_\eta^2} \right) \cdot \left[ 1 - \tanh\left( \beta_2 \, \kappa_f \right) \right]^{-1}.
\end{equation}
Here, $\beta_1$ and $\beta_2$ are campaign-specific empirical constants, $\sigma_\eta$ represents a characteristic modal variance scale, and $\kappa_f = |d^2\eta_3/d\eta_1^2|$ is the fold-proximal curvature indicator.

Physically, when the boundary layer resides on the flat, resilient portions of the coupled sheet, $\kappa_f \to 0$ and $\eta_3$ variations remain small, reducing $\Phi \to 1$ and returning control to Yamada's local gradient physics. However, as the trajectory approaches a critical topological cliff, the rapid inflation of global curvature $\eta_3$ and the sharp spikes in fold proximity $\kappa_f$ mathematically scale the stability functions upward. This non-local regularization dynamically sustains a baseline level of turbulent exchange proportional to the integrated stress accumulation of the entire column. By using the unfolded manifold coordinates to supervise the local closure, the parameterized atmospheric boundary layer is prevented from undergoing numerical flickering or premature collapse, preserving the continuous, intermittent orbital dynamics observed across the field campaigns.

\subsection{Translation table: Geometric properties and atmospheric manifestations}

\begin{tabular}{@{}ll@{}}
Folded manifold & $\to$ Multiple atmospheric equilibria \\
Fold proximity & $\to$ Reduced resilience and heightened intermittency \\
Branch jump & $\to$ Rapid regime transition \\
Hysteresis loop & $\to$ Path-dependent atmospheric evolution \\
Curvature growth ($\eta_3$ inflation) & $\to$ Structural stress approach to turbulence collapse \\
Brittle fold edge & $\to$ Abrupt decoupling over flat terrain \\
Rubbery fold edge & $\to$ Lingering intermittency over complex terrain \\
\end{tabular}

\subsection{Intermittency as orbital dynamics: A provisional geometric perspective}

\textit{This subsection presents a speculative geometric interpretation of boundary-layer intermittency based on phase-space circulation patterns. While the orbital framework is mathematically consistent with the observed manifold topology, full validation of deterministic limit-cycle properties requires multi-decade statistical analysis across all campaigns and is deferred to future work.}

The folded manifold geometry suggests a reframing of classical intermittency phenomena. Rather than pure stochastic ``collapse-and-restart'' dynamics, intermittent bursting can be understood as deterministic, recurring excursions within the state space, governed by alternating mechanical and radiative control loops. This orbital interpretation is visualized on the phase portrait of Figure~\ref{fig:phase-portraits} via circulation annotations:

\begin{enumerate}
  \item \textbf{Radiative Control Phase:} Intense surface cooling weakens turbulent transport, driving the trajectory away from the well-mixed region and onto the decoupled equilibrium sheet. Local gradient metrics saturate as stratification increases, yet the manifold coordinates continue to evolve.
  \item \textbf{Curvature Accumulation:} Frictional isolation allows low-level jet acceleration aloft, sharpening vertical gradients and inflating the curvature mode $\eta_3$ while $\Rig$ remains blind.
  \item \textbf{Mechanical Reconnection:} Accumulated vertical shear eventually triggers a transient snapback toward fold-proximal regions, generating episodic shear-burst reconnection. This transient restores surface-to-air connectivity before radiative control resumes dominance.
\end{enumerate}

This circulation description is offered as a diagnostic lens for interpreting observed intermittency patterns. However, the determination of whether this orbit closes as a formal limit cycle, or whether it represents merely a recurring geometric pattern in aperiodic trajectories, requires extensive statistical validation across the full multi-campaign dataset and is marked for future investigation.

\section{Computational Framework as an Evidence Engine}

\subsection{Pipeline overview and strategic positioning}

The SpectralBL-Analytics pipeline executes as a five-stage evidence cascade, transforming raw heterogeneous field measurements into low-dimensional dynamical system equations, thereby bridging data collection and invariant manifold tracking. Each stage is designed to preserve physical information fidelity while progressively extracting dynamical structure. The pipeline has been validated across four campaigns (CASES-99, FLOSS, GABLS3, BLLAST) and optimized for deployment in production atmospheric analysis workflows.

\subsection{Stage 1: Campaign Ingestion and Profile Reconstruction}

Raw measurements from disparate instrument networks (tower anemometers, thermistors, sonic anemometers, model outputs) are ingested via campaign-specific parsers, interpolated onto a unified p-FEM grid, and reconstructed as smooth, physically consistent vertical profiles. This stage eliminates sensor clustering artifacts through the p-FEM mass-matrix formalism described earlier. The CASES-99 campaign, for example, ingests 6,538 total samples spanning 685.9 hours of nocturnal evolution, distributed across seven discrete tower heights from 1.5\,m to 50\,m. The piecewise linear interpolation to the p-FEM basis ensures that both coarse-resolution campaigns (BLLAST, sparse tower networks) and dense model-level simulations (GABLS3 with 70+ vertical levels) are mapped into a common functional space. The resulting reconstructed profiles are smooth to machine precision and admit well-defined derivatives.

\subsection{Stage 2: Spectral Decomposition and Manifold Coordinate Extraction}

The Singular Value Decomposition (SVD), constrained by the p-FEM mass matrix $\mat{M}$, extracts leading modes $\eta_1(\vect{x}, t)$, $\eta_2(\vect{x}, t)$, $\eta_3(\vect{x}, t)$ representing integrated column structure. For CASES-99, this yields an effective dimension $D_{\text{eff}} = 2.0537$, demonstrating extremely low dimensionality. The Shannon entropy of the singular-value distribution is $H = 0.2110\,\mathrm{bits}$, confirming that the manifold geometry is tightly constrained by the coupled feedback between surface cooling and frictional dynamics. The condition number of the reconstructed correlation matrix is $\kappa \approx 6.85$, well-conditioned and free from numerical pathologies. These quantitative diagnostics validate that the observed folded structure is robust and not an artifact of ill-posed regression or overfitting.

\subsection{Stage 3: Trajectory Assembly and Information-Theoretic Diagnosis}

Time-ordered sequences of $(\eta_1, \eta_2, \eta_3)$ are assembled to form continuous campaign trajectories in the reduced state space. Mutual information (MI) between the curvature mode $\eta_3$ and low-level Richardson proxies is computed across sliding windows to quantify predictive structure. In BLLAST, the optimal embedding delay is $\tau_{\text{best}} \approx 0.0\,\mathrm{h}$ with lag correlation $R_\tau \approx 0.248$ and mutual information $I(\eta_3; \Rig \leq 0.25) \approx 1.453\,\mathrm{bits}$. This high subcritical MI demonstrates that the entire atmospheric column undergoes coherent structural reconfiguration as it approaches the critical Richardson threshold. The timing synchronicity validates the geometric framework: transitions are not localized, point-wise instabilities, but global state-manifold events where the curvature mode $\eta_3$ acts as a precursor registering column-integrated stress accumulation while local metrics saturate.

\subsection{Stage 4: Sparse Identification of Nonlinear Dynamics (SINDy-BVP)}

The Stage 4 operator discovery phase applies SINDy with boundary-value problem constraints (SINDy-BVP) to automatically extract governing equations from the manifold trajectories. SINDy searches for the sparsest linear combination of candidate nonlinear basis functions (polynomials, trigonometric functions, spatially-varying products) that simultaneously:
\begin{enumerate}
  \item Reproduces the observed time derivatives $d\eta_k/dt$ with minimal residual error.
  \item Respects boundary conditions at the upper and lower domain edges, encoded as hard constraints.
  \item Achieves maximal sparsity---ideally 15--20\% non-zero coefficients in the operator matrix.
\end{enumerate}

The SINDy-BVP solution for CASES-99 achieves 19.5\% sparsity (approximately 80 non-zero coefficients out of 410 basis terms) with all identified coefficients bounded by $|c_i| \leq 0.1836$ to reflect the normalized dynamics of the reduced-coordinate system. This high level of sparsity confirms that nocturnal boundary-layer transitions are organized by a low-order set of coupled physical processes rather than a full-rank, unstructured system. The sparse structure is physically interpretable:
\begin{itemize}
  \item Shear-production terms (proportional to $\eta_1 \eta_2$) dominate $d\eta_1/dt$.
  \item Stratification-feedback terms (proportional to $\eta_3$) regulate $d\eta_2/dt$ and trigger fold-proximal basin transitions.
  \item Cross-coupling terms encode the structural deformation loop by which shear accumulation (tracked in $\eta_3$) and radiative cooling interact.
\end{itemize}

A crucial innovation of the SINDy-BVP approach is the enforcement of spatially-varying operator structure. Rather than assuming a global, height-independent operator, the method allows coefficient tensors to vary across the vertical domain via polynomial basis expansions. This allows the discovered equations to encode the phenomenology that wind-shear-driven production dominates the lower boundary layer while buoyancy-and-curvature-driven redistribution dominates the upper column. The result is a spatially-localized operator map that transitions smoothly from lower-layer dynamics (dominated by mechanical forcing) to upper-layer dynamics (radiatively modulated), exactly mirroring the observed vertical structure of stable boundary-layer transitions.

\subsection{Stage 5: Bifurcation Analytics and Stability Cartography}

The sparse operator equations are continued as families of steady-state solutions, swept across external forcing parameters (geostrophic wind magnitude, surface cooling rate, roughness length). The continuation traces out the multi-valued branch structure, quantifies local Jacobian eigenvalues to detect limit-point transitions, and maps out stability contours in the forcing parameter space. The fold-line structure---the set of points where local bifurcations occur---is extracted geometrically and compared to the observed transition corridors in the raw data. The transversality derivatives characterizing fold geometry (e.g., $\mathcal{T}_{\mathrm{CASES-99}} \approx -0.4110$ vs. $\mathcal{T}_{\mathrm{FLOSS}} \approx -0.0713$) are computed via numerical branch differentiation, quantifying the terrain-dependent stiffness of the fold edge and predicting transition character (brittle versus rubbery).

\subsection{Takens Embedding Theorem and Attractor Reconstruction}

The theoretical foundation for manifold reconstruction from scalar time series rests on Takens' embedding theorem \citep{takens1981}. Given a single time series observable $y(t)$---such as temperature or velocity at a fixed height---the theorem states that under generic conditions, the $d$-dimensional attractor $\mathcal{A}$ can be faithfully reconstructed via delay embedding:
\begin{equation}
\vect{Y}(t) = [y(t), y(t - \tau), y(t - 2\tau), \ldots, y(t - (d_e - 1)\tau)]^\top \in \mathbb{R}^{d_e},
\end{equation}
where $\tau$ is the embedding delay, $d_e$ is the embedding dimension, and provided that $d_e > 2d_A + 1$ (where $d_A$ is the true attractor dimension), the reconstructed manifold $\mathcal{Y}$ is diffeomorphic to $\mathcal{A}$.

In our analysis pipeline, the automated delay-estimation module employs mutual information optimization to search for the delay $\tau$ that balances two competing objectives: sufficient decorrelation between time-lagged samples (avoiding redundancy) and preservation of dynamical coupling (avoiding excessively distant time steps that lose continuity). The DelayEstimator module identifies zero crossings of the autocorrelation function to initialize the search, then refines the estimate through global minimization of the average mutual information over a candidate delay set. For Markovian (stationary) regimes in the SBL, the optimal delay frequently approaches zero, indicating that the column's reduced state can be estimated from instantaneous vertical structure alone, without historical dependence. For history-dependent (intermittent) regimes, non-zero optimal delays emerge, quantifying the memory horizon of the atmospheric column's dynamics.

This embedding framework is critical because it bridges two conceptually distinct views of the data: (1) the high-dimensional, multi-level measurement space of the atmospheric column, and (2) the low-dimensional attractor that organizes temporal evolution. The embedding theorem guarantees that if one has sufficient observational bandwidth (either many heights or sufficient temporal sampling), the topological properties of the true attractor (dimension, folding pattern, branch structure) are faithfully preserved in the reconstructed embedding. Conversely, if observational data are too sparse (too few measurement heights or too coarse temporal resolution), the embedding dimension required to unfold the attractor grows arbitrarily, signaling that the data are insufficient for reliable reconstruction.

The embedding delay-estimation process in DelayEstimator also provides a diagnostic tool for regime classification. The estimated delay $\tau^*$ and its associated mutual information profile can distinguish between:
\begin{itemize}
  \item \textbf{Markovian regimes:} $\tau^* \approx 0$, indicating local balance between vertical gradients and turbulent transport; memory-less evolution.
  \item \textbf{History-dependent regimes:} $\tau^* > 0$, indicating that past profile structure influences current evolution; typical of highly stratified, slowly-adjusting configurations.
  \item \textbf{Intermittent regime boundaries:} $\tau^*$ exhibits rapid variation as the trajectory crosses fold-proximal zones, reflecting the sharp increase in system memory as mechanical and radiative forcing become mutually dependent.
\end{itemize}

\subsection{Integration with the low-dimensional attractor picture}

The p-FEM-based $(\eta_1, \eta_2, \eta_3)$ coordinate system derived in Stage 2 can be understood as a principal-component unfolding of the full attractor manifold $\mathcal{A}$. By operating directly on spatial profiles rather than delay-embedded time series, the p-FEM approach achieves dimensionality reduction without sacrificing spatial information. The reconstructed coordinates automatically encode:
\begin{itemize}
  \item $\eta_1$: Integrated column depth and shear accumulation (analogous to first delay-embedding mode).
  \item $\eta_2$: Residual vertical buoyancy structure (analogous to second delay-embedding mode).
  \item $\eta_3$: Higher-order curvature and profile distortion (analogous to third or fourth embedding mode).
\end{itemize}

This geometric interpretation validates the use of a three-dimensional reduced coordinate system: Takens' theorem predicts that $d_e > 2 d_A + 1$, so if $d_A \approx 1.5$ (consistent with the effective dimension $D_{\text{eff}} = 2.0537$ from CASES-99), then $d_e > 4$ is recommended. In practice, employing $d_e = 3$ captures the dominant branching structure and fold geometry, with the understanding that sub-dominant transverse modes (Lyapunov-exponent-ordered) contribute perturbations that remain small for the regime-transition predictions of interest. Extensions to $d_e = 4$ or $d_e = 5$ are straightforward using the same computational pipeline and would further refine the characterization of intermittent bursting.

\subsection{p-FEM mass matrix and functional geometry}

The p-FEM mass matrix $\mat{M}$ is constructed via numerical integration over the polynomial basis elements:
\begin{equation}
M_{ij} = \int_{z_{\text{min}}}^{z_{\text{max}}} \phi_i(z) \phi_j(z) \, \rho(z) \, dz,
\end{equation}
where $\phi_k$ are Chebyshev polynomials and $\rho(z)$ is a weighting function (typically taken as uniform, but can incorporate terrain-dependent density profiles). This matrix encodes the continuous Hilbert-space inner product between basis functions. By performing SVD in the $\mat{M}$-weighted metric, the resulting singular vectors $\vect{v}_k$ satisfy $\vect{v}_i^\top \mat{M} \vect{v}_j = \delta_{ij}$, meaning they are orthonormal in functional space rather than Euclidean space. This invariance to sensor-height clustering is the key technical innovation that allows CASES-99 (7 discrete tower heights) and GABLS3 (70+ model levels) to be analyzed within the same unified coordinate system without spurious metric distortions.

\section{Risk Assessment and Methodological Limitations}

\subsection{Noise sensitivity and measurement fidelity}

Field measurements contain inherent instrumental noise, sensor drift, and calibration uncertainties. The p-FEM reconstruction process acts as a polynomial smoother, which can suppress high-frequency noise but may also attenuate genuine sharp gradients if the polynomial basis degree is too low. Conversely, excessively high polynomial orders (degree $> 8$) can induce spurious oscillations (Runge phenomenon) near measurement boundaries. The current implementation uses degree-6 Chebyshev bases, a balance validated empirically across all four campaigns. Future sensitivity analysis should quantify how manifold topology changes under synthetic noise perturbations (e.g., $\pm 5\%$ random noise added to original profiles) to establish confidence bounds on transition estimates.

\subsection{Boundary layer closure and inner-boundary handling}

The log-law hypothesis ($u(z) = u_*/\kappa \ln(z/z_0)$ in the inner layer) is enforced through hard Dirichlet boundary conditions at $z = z_0$ (roughness length). When $z_0$ is misspecified (either through instrumental error or true spatial variability at the footprint scale), the entire p-FEM reconstruction is shifted systematically. This can affect the absolute Richardson values and the precise location of the fold edge in external forcing space. For CASES-99, $z_0 \approx 0.02$\,m is well-documented, but for heterogeneous terrain (FLOSS, BLLAST), $z_0$ may vary by factor of 2-3 across the measurement footprint. A next-generation analysis should employ adaptive $z_0$ estimation coupled with spatial aggregation strategies to mitigate this systematic error source.

\subsection{Data sufficiency and embedding dimension}

Takens' theorem requires an embedding dimension $d_e > 2 d_A + 1$. In CASES-99, $d_A \approx 1.5$ suggests $d_e > 4$, yet the analysis uses $d_e = 3$. While the three-dimensional folded structure captures the dominant bifurcation geometry, subordinate transverse modes (smaller-amplitude Lyapunov exponents) are unresolved. These missing dimensions contribute to the residual scatter observed around the manifold branches and may explain the finite width of the fold-proximal transition corridors. A rerun with $d_e = 4$ or $d_e = 5$ is recommended for future work, contingent on whether higher-dimensional branches represent genuine dynamical structure or merely noise-aliased echoes.

\subsection{Collinearity in SINDy basis and operator uniqueness}

The SINDy-BVP stage discovers equations by solving a sparse least-squares problem over candidate basis functions. When basis functions are correlated (multicollinearity), multiple sparse solutions can fit the data approximately equally well, leading to operator non-uniqueness. This is particularly acute in polynomial bases, where low-degree terms (e.g., $\eta_1^2$ and $\eta_1 \cdot \eta_2 / \eta_1$) can be nearly linearly dependent. The current implementation employs orthogonalized basis and iterative thresholding to mitigate this; however, a formal analysis of the solution's stability under small perturbations to input trajectories has not been performed. If operator coefficients exhibit large sensitivity to minute training-data changes, confidence in the extrapolated bifurcation predictions decreases. A sensitivity analysis using bootstrap resampling of the CASES-99 trajectory should quantify this uncertainty.

\subsection{Generalization across campaigns: physics-agnostic versus physics-aware learning}

The SINDy-BVP operators discovered from CASES-99 do not directly transfer to FLOSS or BLLAST. This is by design: the coefficient structure is physics-aware, encoding terrain-dependent effects and domain-heterogeneous forcing. However, the manifold topology (fold structure, branch organization, transversality) is robust across campaigns. This distinction between global geometry (preserved) and local operator coefficients (campaign-specific) suggests that the principal predictive target---regime transition timing and character---lies in the geometry rather than precise operator values. For future operational applications (e.g., real-time boundary-layer forecast guidance), the geometric framework offers a portable diagnostic, while campaign-specific operators should be retrained for local deployment.

\subsection{Intermittency model validation}

The orbital dynamics interpretation of intermittency (Section 4.5) frames bursting as deterministic limit-cycle excursions rather than stochastic collapse-and-restart. This hypothesis is plausible but remains partially validational: full proof would require identifying closed orbits in the phase space and demonstrating their stability over multi-year statistical ensembles. The BLLAST campaign provides only a single season of observations (late afternoon/evening decay), insufficient for rigorous periodicity analysis. Multi-decade records or controlled laboratory experiments (wind-tunnel stably-stratified flows with automated heating cycles) would be needed to definitively establish whether intermittency obeys deterministic or primarily stochastic laws.

\subsection{Generalization to global scales and climate transitions}

The stable boundary layer is a local, surface-proximate phenomenon. Scaling arguments from 1\,m to 100\,m (typical SBL scales) to planetary atmospheric scales ($>1000$\,km) involves fundamental changes in the governing force balances (Coriolis acceleration becomes comparable to pressure-gradient forces, vertical mixing is replaced by large-scale baroclinic circulation). The folded-manifold framework may not extend to synoptic-scale ``stably stratified'' regimes (e.g., polar inversions, upper-tropospheric stabilities). This analysis is deliberately local and campaign-focused; any extrapolation to climate or global models remains speculative.

\subsection{Recommendation for risk mitigation}

To increase confidence in the manifold-based transition framework prior to operational deployment, we recommend:
\begin{enumerate}
  \item Conduct synthetic data validation: Add realistic instrumental noise ($\sim 5\%$) to original profiles and verify that manifold topology is preserved.
  \item Perform cross-validation: Train SINDy-BVP operators on half of CASES-99, predict bifurcation structure, compare to held-out second half.
  \item Extend embedding analysis to $d_e = 4, 5$ to test whether higher-dimensional geometry resolves sub-dominant intermittency modes.
  \item Apply independent event-detection algorithms (wavelet-based intermittency detection, event-counting statistics) and cross-validate against geometric regime boundaries.
  \item Conduct controlled laboratory experiments (wind-tunnel stably-stratified flows) to test predictions of transition timing and character under perfectly controlled forcing.
\end{enumerate}

\section{Synthesis and Conclusions}

The principal contribution is a physics result unified by geometric reasoning: the nocturnal boundary layer exhibits a folded, low-dimensional state-space geometry that organizes transition, intermittency, and hysteresis across campaigns with vastly different forcing, terrain, and instrumentation contexts. This geometry is not an artifact of measurement or analysis choice; it persists across the spectrum from grassland (CASES-99) to alpine snowpack (FLOSS) to idealized simulation (GABLS3) to rapidly-forced decay scenarios (BLLAST).

In the classical framework, Richardson diagnostics are viewed as fundamental---local gradient stability metrics that trigger or suppress turbulence through pointwise closure laws. When $\Rig$ exceeds a critical threshold $\Ric$, local turbulence is formally unstable and must collapse. In the geometric framework presented here, Richardson diagnostics remain useful but incomplete. They act as projection operators onto a lower-dimensional diagnostic plane, and near fold singularities, this projection induces information loss: multiple distinct atmospheric column states project onto nearly identical Richardson values. The familiar S-curve is therefore not a property of the atmosphere---it is a property of the projection. Equivalently, classical Richardson thresholds are not rejected; they are reinterpreted as local manifestations of global manifold structure.

The strategic evidence base supporting this reframing comprises five stages:

\begin{enumerate}
  \item \textbf{Stage 1 (Ingestion)} reconstructs heterogeneous measurements into smooth functional spaces, eliminating sensor-clustering artifacts through p-FEM interpolation and mass-matrix weighting.
  \item \textbf{Stage 2 (Spectral Decomposition)} extracts low-dimensional coordinates $(\eta_1, \eta_2, \eta_3)$ via SVD in functional metric space, yielding CASES-99 effective dimension $D_{\text{eff}} = 2.0537$ and entropy $H = 0.2110$\,bits---extraordinarily low values confirming tight phase-space organization.
  \item \textbf{Stage 3 (Information-Theoretic Diagnosis)} reveals high subcritical mutual information ($I \approx 1.453$\,bits between $\eta_3$ and $\Rig$ in BLLAST) and near-zero embedding delay ($\tau_{\text{best}} \approx 0.0$\,h), proving that atmospheric column reconfiguration is both coherent and Markovian.
  \item \textbf{Stage 4 (SINDy-BVP)} automatically discovers governing equations achieving 19.5\% sparsity (80 non-zero coefficients out of 410) with spatially-varying operator structure, confirming that regime transitions emerge from interpretable low-order physical processes: shear production, buoyancy feedback, and curvature-driven redistribution.
  \item \textbf{Stage 5 (Bifurcation Analytics)} traces multi-valued branch structure, characterizes fold geometry via transversality derivatives ($\mathcal{T}_{\text{CASES-99}} \approx -0.4110$ vs. $\mathcal{T}_{\text{FLOSS}} \approx -0.0713$), and maps terrain-dependent modulation of transition character from brittle (rapid decoupling) to rubbery (lingering intermittency).
\end{enumerate}

This five-stage pipeline functions as an evidence engine, transforming raw disparate measurements into validated physical insights. The geometric interpretation resolves long-standing tensions in stable-boundary-layer science:

\begin{itemize}
  \item \textbf{Why do nearly identical forcings yield disparate transitions?} Because terrain heterogeneity modulates the manifold curvature (transversality), creating brittle versus rubbery fold geometries with fundamentally different dynamical signatures.
  \item \textbf{Why do local Richardson thresholds fail to predict transitions in strongly stable regimes?} Because the manifold folds, and high-dimensional columns project onto nearly identical Richardson values, saturating the diagnostic.
  \item \textbf{Why is intermittency structured rather than random?} Because the state trajectory executes deterministic excursions within a constrained phase space, recurrently visiting fold-proximal regions where mechanical and radiative controls alternate.
\end{itemize}

By unifying McNider's pioneering bifurcation models, van de Wiel's maximum sustainable heat flux framework, and modern embedding theorems (Takens, SINDy, continuation methods), the present work reframes regime transitions from arbitrary threshold crossings into deterministic traversals of a universal, empirically-validated manifold surface. This geometric view is not merely a mathematical convenience; it yields immediate practical dividends: predictive power for transition timing and character, transportable diagnostics across diverse field campaigns, and a natural foundation for next-generation boundary-layer parameterizations that interpolate between local and global constraints.

Future work should extend the analysis to the full multi-campaign ensemble, validate bifurcation predictions against independent event-detection schemes, and test whether the geometric framework extends to other stably-stratified regimes (polar inversions, upper-tropospheric stabilities) or remains specific to the surface nocturnal boundary layer. Laboratory experiments with controlled forcing profiles and multi-year field campaigns with near-identical instrument networks would further strengthen the evidence base. The manifold geometry offers a physics-first, data-validated foundation for this next generation of atmospheric boundary-layer understanding.

\bibliography{refs}
\bibliographystyle{ametsoc2014}

\end{document}